% ===============================================================
%  جلسه نوزدهم
% ===============================================================
\chapter{آمار: مفاهیم تجربی، تخمین، آزمون فرض و پیش‌بینی (تعمیق)}

\begin{note}
\notetitle اسلایدهای جلسه نوزدهم همان مباحث بخش «آمار» جلسه هجدهم (آمار در مقابل احتمال، توزیع آماری در مقابل تجربی، سرفصل‌های اصلی آمار و بهینگی) را مرور و تکمیل می‌کنند. برای پرهیز از تکرار، در این فصل خلاصه‌ای از آن مطالب آورده و سپس هر سرفصل را با جزئیات بیشتر، مثال‌های حل‌شده و ابزارهای لازم برای جلسات بعد (تخمین نقطه‌ای و تخمین میانگین مربع) گسترش می‌دهیم.
\end{note}

\section{خلاصه: آمار در مقابل احتمال}
\begin{itemize}
\item \textbf{احتمال} علمی مجرد است: با اصول موضوعه (فضای احتمال $(\Omega,\F,P)$ و اصول کولموگروف) شروع می‌کند، تعاریف جانبی (احتمال شرطی، استقلال، متغیر تصادفی، \ldots) می‌سازد و احکامی (قانون بیز، قانون زنجیره‌ای، قضیه حد مرکزی، قانون اعداد بزرگ، \ldots) را اثبات می‌کند.
\item \textbf{آمار} کاربرد احتمال در دنیای واقع است، جایی که مدل دقیق را نداریم و ناچار به «استدلال‌های استقرایی بر پایه مشاهدات» هستیم. آمار در حالت کلی علم کار کردن با داده‌های واقعی است.
\end{itemize}

\begin{center}
\begin{tikzpicture}[>=Stealth,box/.style={draw,very thick,rounded corners,minimum width=3.4cm,minimum height=1.2cm,align=center}]
\node[box,fill=blue!8] (M) at (0,0) {مدل احتمالاتی\\ $f_X(x;\theta)$};
\node[box,fill=orange!10] (D) at (8,0) {داده‌ها\\ $X_1,\ldots,X_n$};
\draw[->,very thick,blue!70!black] (M.north east) to[bend left=20] node[above]{احتمال (قیاسی)} (D.north west);
\draw[->,very thick,orange!70!black] (D.south west) to[bend left=20] node[below]{آمار (استقرایی)} (M.south east);
\end{tikzpicture}
\end{center}

\section{مفاهیم تجربی و ویژگی‌های آن‌ها}

\subsection{توزیع انباشته تجربی}
اگر $X_1,\ldots,X_n$ نمونه‌های \lr{i.i.d.} از $X$ با \lr{CDF} آماری $F_X$ باشند، توزیع انباشته تجربی و میانگین تجربی عبارت‌اند از:
\[
\hat F_X(x)=\frac1n\sum_{i=1}^{n}\ind(X_i\le x),\qquad \bar X_n=\frac1n\sum_{i=1}^{n}X_i .
\]
$\hat F_X$ یک تابع پله‌ای است که در هر نمونه $X_i$ به اندازه $\frac1n$ بالا می‌رود؛ یعنی توزیعی است که به هر نمونه مشاهده‌شده جرم احتمال $\frac1n$ می‌دهد. پس $\bar X_n$ دقیقاً میانگین همین توزیع تجربی است.

\begin{theorem}[ویژگی‌های \lr{CDF} تجربی]
برای هر $x$ ثابت، $n\hat F_X(x)\sim\mathrm{Bin}\big(n,F_X(x)\big)$. در نتیجه:
\[
\E\big[\hat F_X(x)\big]=F_X(x),\qquad \Var\big[\hat F_X(x)\big]=\frac{F_X(x)\big(1-F_X(x)\big)}{n}\xrightarrow{n\to\infty}0 .
\]
\end{theorem}
\begin{proof}
$Y_i=\ind(X_i\le x)$ متغیرهای برنولی \lr{i.i.d.} با احتمال موفقیت $\Prob(X_i\le x)=F_X(x)$ هستند. $n\hat F_X(x)=\sum Y_i$ جمع $n$ برنولی مستقل، یعنی دوجمله‌ای است. میانگین و واریانس از جدول جلسه دهم به دست می‌آیند.
\end{proof}

\begin{note}
\notetitle این قضیه پاسخ پرسش «\textbf{چرا} مفاهیم تجربی متغیر تصادفی‌اند؟» را دقیق‌تر می‌کند: $\hat F_X(x)$ برای هر $x$ یک متغیر تصادفی (دوجمله‌ای تقسیم بر $n$) است. «به طور متوسط درست» است (میانگینش برابر مقدار آماری است) و پراکندگی‌اش با $\frac1{\sqrt n}$ کم می‌شود. بیشترین پراکندگی در میانه توزیع ($F_X(x)=\frac12$) رخ می‌دهد.
\end{note}

\subsection{هیستوگرام: تقریب تجربی چگالی}
برای تقریب \lr{PDF}، محور را به بازه‌های $[a_k,a_k+h)$ تقسیم و کسر نمونه‌های هر بازه را بر $h$ تقسیم می‌کنیم:
\[
\hat f_X(x)=\frac{1}{nh}\#\{i:\ X_i\in[a_k,a_k+h)\},\qquad x\in[a_k,a_k+h).
\]
طبق قانون اعداد بزرگ، وقتی $n\to\infty$ این مقدار به $\frac{1}{h}\Prob(X\in[a_k,a_k+h))$ میل می‌کند که برای $h$ کوچک تقریباً $f_X(x)$ است (تعریف چگالی در جلسه هفتم). دوباره یک «مصالحه» داریم: $h$ کوچک دقت مکانی بیشتر ولی نوسان تصادفی بیشتر (نمونه‌های کمتر در هر بازه) می‌دهد.

\subsection{واریانس تجربی و ضریب $\frac{1}{n-1}$}
\begin{theorem}[واریانس نمونه‌ای]
برای نمونه‌های \lr{i.i.d.} با واریانس $\sigma^2$:
\[
\E\Big[\frac1n\sum_{i=1}^n(X_i-\bar X_n)^2\Big]=\frac{n-1}{n}\sigma^2,\qquad\text{پس}\qquad
S_n^2\triangleq\frac{1}{n-1}\sum_{i=1}^n(X_i-\bar X_n)^2\ \Rightarrow\ \E\big[S_n^2\big]=\sigma^2 .
\]
\end{theorem}
\begin{proof}
با اتحاد $\sum(X_i-\bar X)^2=\sum(X_i-\mu)^2-n(\bar X-\mu)^2$ (بسط دهید و از $\sum(X_i-\mu)=n(\bar X-\mu)$ استفاده کنید):
\[
\E\sum(X_i-\bar X)^2=n\sigma^2-n\Var(\bar X)=n\sigma^2-n\cdot\frac{\sigma^2}{n}=(n-1)\sigma^2 .
\]
\end{proof}
\textbf{شهود:} انحراف‌ها از $\bar X_n$ (که خودش از روی همین داده‌ها و «نزدیک به آن‌ها» انتخاب شده) به طور سیستماتیک کوچک‌تر از انحراف‌ها از $\mu$ واقعی‌اند؛ تقسیم بر $n-1$ این کم‌برآوردی را جبران می‌کند. (به $n-1$ «درجات آزادی» می‌گویند: یک درجه آزادی صرف تخمین میانگین شده است.) این مفهوم در جلسه بعد با نام «گرایش» (\lr{Bias}) دقیق می‌شود.

\section{تخمین \lr{(Estimation)}}

متغیر تصادفی $X\sim f_X(x;\theta)$ با پارامتر(های) نامعلوم $\theta$ و نمونه‌های مستقل $X_1,\ldots,X_n$ را در نظر بگیرید.

\subsection{تخمین نقطه‌ای}
هدف: تابعی $\hat\theta=\hat\theta(X_1,\ldots,X_n)$ که «به‌خوبی» $\theta$ را تقریب بزند. مثال‌ها:
\begin{center}
\renewcommand{\arraystretch}{1.4}
\begin{tabular}{l|l|l}
مدل & پارامتر & تخمین‌گر طبیعی\\\hline
$\mathrm{Bern}(p)$ & $p$ & $\hat p=\bar X_n$ (فرکانس نسبی)\\
$\mathcal N(\mu,\sigma^2)$ & $\mu,\ \sigma^2$ & $\hat\mu=\bar X_n,\quad\hat\sigma^2=S_n^2$\\
$\mathrm{Exp}(\lambda)$ & $\lambda$ & $\hat\lambda=1/\bar X_n$ (چون $\E X=1/\lambda$)\\
$\mathcal U(0,\theta)$ & $\theta$ & $\hat\theta=\max_iX_i$ \ یا\ $\hat\theta=2\bar X_n$\\
$\mathrm{Pois}(\lambda)$ & $\lambda$ & $\hat\lambda=\bar X_n$
\end{tabular}
\end{center}
دقت کنید که برای یک پارامتر می‌توان چند تخمین‌گر متفاوت ساخت (مثل $\mathcal U(0,\theta)$)؛ اینکه کدام «بهتر» است، به معیار بهینگی بستگی دارد (بخش آخر و جلسه بعد).

\subsection{تخمین بازه‌ای}
هدف: بازه تصادفی $(\theta_l,\theta_u)$ با $\Prob(\theta_l<\theta<\theta_u)\ge1-\delta$.

\begin{theorem}[بازه اطمینان برای میانگین]
اگر $X_i\sim\mathcal N(\mu,\sigma^2)$ با $\sigma$ معلوم باشند (یا به طور تقریبی، طبق \lr{CLT}، برای $n$ بزرگ و هر توزیعی):
\[
\Prob\Big(\bar X_n-z_{\delta/2}\frac{\sigma}{\sqrt n}<\mu<\bar X_n+z_{\delta/2}\frac{\sigma}{\sqrt n}\Big)=1-\delta,
\]
که در آن $z_{\delta/2}$ عددی است که $\Prob(Z>z_{\delta/2})=\frac\delta2$ (مثلاً $z_{0.025}=1.96$ برای ۹۵٪ و $z_{0.005}=2.576$ برای ۹۹٪).
\end{theorem}
\begin{proof}
$Z=\frac{\bar X_n-\mu}{\sigma/\sqrt n}\sim\mathcal N(0,1)$. پس $\Prob(|Z|<z_{\delta/2})=1-\delta$ و با جابه‌جا کردن نامساوی برای $\mu$ حکم به دست می‌آید.
\end{proof}

\begin{warning}
\warntitle تعبیر درست بازه اطمینان: $\mu$ عددی ثابت (غیرتصادفی) است و این \emph{بازه} است که تصادفی است. «بازه ۹۵٪» یعنی اگر آزمایش نمونه‌گیری را بارها تکرار کنیم، حدود ۹۵٪ بازه‌های ساخته‌شده $\mu$ را دربر می‌گیرند؛ نه اینکه «$\mu$ با احتمال ۹۵٪ در این بازه مشخص است».
\end{warning}

\begin{example}[نظرسنجی]
در یک نظرسنجی از ۱۶۰۰ نفر، ۸۸۰ نفر موافق بودند. بازه اطمینان ۹۵٪ برای نسبت موافقان در جامعه را بیابید.
\begin{solution}
$\hat p=\frac{880}{1600}=0.55$. انحراف معیار $\hat p$ برابر $\sqrt{p(1-p)/n}\approx\sqrt{0.55\times0.45/1600}\approx0.0124$. بازه: $0.55\pm1.96\times0.0124=0.55\pm0.024$، یعنی $[0.526,\ 0.574]$. («حاشیه خطای» نظرسنجی حدود $2.4\%$ است.)
\end{solution}
\end{example}

\subsection{تخمین بیزی}
در دیدگاه بیزی، $\theta$ خودش متغیر تصادفی با توزیع پیشین $f_\Theta$ است و پس از مشاهده داده‌ها توزیع پسین را حساب می‌کنیم:
\[
f_{\Theta|X_{1:n}}(\theta\given x_{1:n})=\frac{f_\Theta(\theta)\prod_{i}f_X(x_i;\theta)}{\int f_\Theta(t)\prod_if_X(x_i;t)\dd t}.
\]
از روی پسین می‌توان تخمین نقطه‌ای (میانگین پسین، مُد پسین یا میانه پسین) و «بازه باورپذیری» ساخت.

\begin{example}[پسین بتا و میانگین پسین]
برای سکه‌ای با احتمال شیر نامعلوم و پیشین یکنواخت، پس از $k$ شیر در $n$ پرتاب، پسین $\mathrm{Beta}(k+1,n-k+1)$ است (جلسه نهم). میانگین پسین $\frac{k+1}{n+2}$ و مُد پسین $\frac{k}{n}$ است. برای $n=3,k=3$: تخمین فرکانسی $\hat p=1$ (یعنی «سکه هرگز خط نمی‌آورد»!)، اما میانگین پسین $\frac45=0.8$ که معقول‌تر است. پیشین مانند «نمونه‌های مجازی» عمل می‌کند و تخمین را در داده‌های کم پایدار می‌سازد.
\end{example}

\section{آزمون فرض \lr{(Hypothesis Testing)}}

\begin{definition}[اجزای آزمون فرض]
\begin{itemize}
\item \textbf{فرض صفر} $H_0$ (مثلاً $\theta=\theta_0$) و \textbf{فرض مقابل} $H_1$ (مثلاً $\theta\ne\theta_0$).
\item یک \textbf{آماره آزمون} $T(X_1,\ldots,X_n)$ و یک \textbf{ناحیه رد} $R$: اگر $T\in R$، $H_0$ را رد می‌کنیم.
\item \textbf{خطای نوع اول:} رد $H_0$ در حالی که درست است؛ احتمال آن \textbf{سطح معناداری} $\alpha$ نام دارد.
\item \textbf{خطای نوع دوم:} نپذیرفتن $H_1$ در حالی که درست است؛ احتمال آن $\beta$ و $1-\beta$ «توان آزمون» است.
\end{itemize}
\end{definition}

\begin{center}
\renewcommand{\arraystretch}{1.4}
\begin{tabular}{c|c|c}
 & $H_0$ درست است & $H_1$ درست است\\\hline
$H_0$ را رد نمی‌کنیم & تصمیم درست ($1-\alpha$) & خطای نوع دوم ($\beta$)\\
$H_0$ را رد می‌کنیم & خطای نوع اول ($\alpha$) & تصمیم درست ($1-\beta$)
\end{tabular}
\end{center}

\begin{note}
\notetitle این همان ساختار مثال «آزمایش پزشکی» جلسه چهارم است: مثبت کاذب (خطای نوع اول) و منفی کاذب (خطای نوع دوم). کاهش هم‌زمان هر دو خطا با $n$ ثابت ممکن نیست؛ با تنگ کردن ناحیه رد، $\alpha$ کم و $\beta$ زیاد می‌شود. تنها راه کاهش هر دو، افزایش تعداد نمونه است.
\end{note}

\begin{example}[آزمون میانگین]
یک کارخانه ادعا می‌کند وزن متوسط بسته‌هایش $\mu_0=500$ گرم با انحراف معیار ۱۰ گرم است. ۲۵ بسته تصادفی وزن متوسط ۴۹۵ گرم دارند. در سطح $\alpha=0.05$ آیا ادعا رد می‌شود؟
\begin{solution}
تحت $H_0$: $Z=\frac{\bar X-500}{10/\sqrt{25}}=\frac{-5}{2}=-2.5$. ناحیه رد (دوطرفه) $|Z|>1.96$ است. چون $|{-2.5}|>1.96$، $H_0$ رد می‌شود. مقدار $p$: $\Prob(|Z|\ge2.5)=2(1-\Phi(2.5))\approx0.012<0.05$.
\end{solution}
\end{example}

\section{پیش‌بینی \lr{(Prediction)}}
در پیش‌بینی، مدل $f_X$ معلوم است و درباره داده‌های آینده سؤال می‌کنیم. دو مثال اسلاید را حل می‌کنیم:

\begin{example}[پیش‌بینی میانگین تجربی]
اگر $n$ نمونه \lr{i.i.d.} از $f_X$ بگیریم، $\bar X_n$ چه توزیعی دارد؟
\begin{solution}
در حالت کلی $\E\bar X_n=\mu$ و $\Var\bar X_n=\frac{\sigma^2}{n}$؛ برای $n$ بزرگ طبق \lr{CLT}: $\bar X_n\approx\mathcal N(\mu,\frac{\sigma^2}{n})$. اگر $f_X$ گاوسی باشد این رابطه دقیق است؛ اگر $f_X$ نمایی با نرخ $\lambda$ باشد، $n\bar X_n$ دقیقاً ارلانگ مرتبه $n$ است (جلسه شانزدهم).
\end{solution}
\end{example}

\begin{example}[پیش‌بینی تعداد موفقیت‌ها]
اگر احتمال موفقیت یک آزمایش $A$ برابر $p$ باشد، تعداد موفقیت‌ها در $n$ تکرار مستقل آینده چه توزیعی دارد؟ برای $p=0.9$ و $n=20$، احتمال اینکه دست‌کم ۱۸ آزمایش موفق باشند چقدر است؟
\begin{solution}
$N\sim\mathrm{Bin}(n,p)$. $\Prob(N\ge18)=\sum_{k=18}^{20}\binom{20}{k}0.9^k0.1^{20-k}=0.2852+0.2702+0.1216\approx0.677$.
\end{solution}
\end{example}

\section{بهینگی: یادآوری}
«بهترین» تخمین‌گر بدون مشخص کردن یک معیار $\ell(\theta,\hat\theta)$ معنا ندارد. سه معیار رایج:
\begin{center}
\renewcommand{\arraystretch}{1.4}
\begin{tabular}{c|c|c}
معیار & فرمول & بهترین ثابت $c$ برای تقریب $X$\\\hline
خطای مربعی & $\E(X-c)^2$ & میانگین $\E X$\\
خطای مطلق & $\E|X-c|$ & میانه\\
خطای صفر--یک & $\Prob(X\neq c)$ & مُد
\end{tabular}
\end{center}
معیار خطای مربعی (\lr{MSE}) به دلیل سادگی جبری و ارتباط مستقیم با واریانس، پرکاربردترین است و موضوع اصلی جلسه بعد خواهد بود.

\section{تمرین‌های تکمیلی}

\begin{exercise}
داده‌های $\{2.1,\ 3.4,\ 1.7,\ 2.9,\ 3.9\}$ از توزیعی نامعلوم به دست آمده‌اند. $\bar X$، $S^2$ و $\hat F_X(3)$ را حساب کنید.
\end{exercise}
\begin{solution}
$\bar X=\frac{14}{5}=2.8$. انحراف‌ها: $-0.7,0.6,-1.1,0.1,1.1$؛ مجموع مربعات $0.49+0.36+1.21+0.01+1.21=3.28$. $S^2=\frac{3.28}{4}=0.82$. $\hat F_X(3)=\frac{3}{5}=0.6$ (سه نمونه $\le3$).
\end{solution}

\begin{exercise}
برای $\mathcal U(0,\theta)$، میانگین تخمین‌گرهای $\hat\theta_1=2\bar X_n$ و $\hat\theta_2=\max_iX_i$ را بیابید. کدام «به طور متوسط» درست است؟
\end{exercise}
\begin{solution}
$\E\hat\theta_1=2\cdot\frac\theta2=\theta$. برای بیشینه، طبق جلسه هفدهم با مقیاس $\theta$: $\E\max=\frac{n}{n+1}\theta$. پس $\hat\theta_1$ به طور متوسط درست است ولی $\hat\theta_2$ همیشه کمی کمتر از $\theta$ تخمین می‌زند (که البته با ضریب $\frac{n+1}{n}$ قابل اصلاح است). مقایسه کامل این دو با معیار \lr{MSE} در جلسه بعد آمده است.
\end{solution}

\begin{exercise}
چند نفر باید در یک نظرسنجی شرکت کنند تا حاشیه خطای ۹۵٪ حداکثر ۲٪ باشد (بدترین حالت $p=0.5$)؟
\end{exercise}
\begin{solution}
$1.96\sqrt{\frac{0.25}{n}}\le0.02\Rightarrow\sqrt n\ge\frac{1.96\times0.5}{0.02}=49\Rightarrow n\ge2401$.
\end{solution}
